Pixel and segment-based change detectors optimize map accuracy
rather than a three-class patch rank with a sentence, and they remain
the right tool for a wall-to-wall disturbance map. Remote-sensing
question answering  and change captioning
 pair imagery with language, but they work on
single-sensor optical imagery and do not check generated text against
measured values. EcoExplain-VLM asks a different question: whether a
general VLM can use the same multi-sensor evidence that a human analyst
would see, and whether its explanation agrees with that evidence.

Remote-sensing VLMs such as GeoChat are trained on
overhead imagery and spatial referring expressions and are a natural
next baseline. They were not evaluated here because the available
compute limited the study to general checkpoints of 0.5--3 billion
parameters. CLIP-style encoders could score the
image pair without generating text, but they would not produce the
justification that the benchmark is designed to test.

The most informative comparison this dataset supports is with simple
baselines. Against the majority class, only Qwen2-VL-2B with the full
card shows a clear margin. Against a single $\Delta$NDVI threshold, none
of the models comes close.

% ============================================================
% 4. CONCLUSIONS
% ============================================================

\section{Conclusions}

EcoExplain-VLM packages 366 Kanha patches, matched 2022/2023 dry-season
Sentinel-2 and Sentinel-1 data, climate and terrain context, and a
three-class NDVI-change label into a reproducible VLM benchmark. On the
55-patch test set, Qwen2-VL-2B reaches 56.4\% zero-shot accuracy with
the full evidence card, and a 16-patch class LoRA raises this to
67.3\%; both figures depend on a card that contains the candidate label.
LLaVA-Interleave-0.5B, PaliGemma-3B and BLIP-2 OPT-2.7B stay at
32.7--36.4\%, the level of always guessing one class, while a single
$\Delta$NDVI threshold reaches 94.5\%. For explanations, free-form
prompts give unusable output, and a structured cite-only prompt makes
explanations consistent with the evidence (75.8\% and 71.1\% correct
claims) at the cost of restating the supplied facts. An explanation LoRA
trained on 16 self-generated justifications gives 61.8\% class accuracy
and 50.5\% correct claims and is a pilot rather than an improvement.

Small open VLMs, as evaluated here, are not a substitute for spectral
change detection in Kanha. Under a constrained prompt they can repeat
the facts they are given. The value of the benchmark is that it makes
this limitation measurable and gives larger or domain-adapted models a
fixed, label-free test against which to show real progress.

% ============================================================
% LIMITATIONS (required by ACL; unnumbered, after the conclusion)
% ============================================================

\section*{Limitations}

\paragraph{Labels.} The classes are percentile ranks of one-year
dry-season NDVI change, not validated disturbance, degradation or
recovery, and the absolute changes are small. Human agreement on 19
patches is 42\%, with a single annotator. An expert audit with a second
annotator is needed before any management use.

\paragraph{Label in the input.} The Condition~C card contains
\texttt{candidate\_class}, and both headline Qwen2-VL-2B results (56.4\%
and 67.3\%) use it. Future evaluations should use the
\texttt{evidence\_only} card by default, which the code now supports.


\paragraph{Sample size and split.} With $n = 55$, 95\% intervals are about
$\pm 13$ percentage points wide. The random stratified split does not
prevent spatial leakage ; a spatial block split should
be added as a second test set.

\paragraph{Models.} The results cover four small open VLMs and say nothing
about larger proprietary or remote-sensing-pretrained models.

\paragraph{Explanation scoring.} Automatic claim matching can miss
paraphrases, can pair a direction word with the wrong indicator, and can
mark a true but small change as incorrect because of the dead-band. The
reported claim correctness is not a human judgement of faithfulness.



\paragraph{Future directions.} Five extensions would change the scientific picture rather than add
rows to a leaderboard: (i)~make label-free evaluation the only primary
table, including Qwen2-VL-2B on all 55 test patches; (ii)~add a spatial
block test; (iii)~collect expert labels for all 366 patches or a
sufficiently large subsample; (iv)~fine-tune on all 256 training
patches without the label line, with early stopping on the development
set; and (v)~evaluate a larger or remote-sensing-pretrained VLM, and a
stronger teacher for explanation distillation, on the same cards.



% ============================================================

% ============================================================
% APPENDIX
% ============================================================

\clearpage
\appendix

\section{Additional Experimental Details}

\subsection{Fine-Tuning Configuration}

Table~\ref{tab:ftconfig} lists the settings shared by the class and
explanation LoRA runs. Both used Qwen/Qwen2-VL-2B-Instruct with the
Hugging Face \texttt{transformers} and \texttt{peft} libraries, and both
trained for one epoch on 16 examples. For the class LoRA, the 16
patches were the configured training size. For the explanation LoRA,
the 16 were the filtered justifications that fell within the first 128
training patches.

\begin{table}[ht]
\centering
\small
\begin{tabular}{@{}ll@{}}
\toprule
Setting & Value \\
\midrule
Base model & Qwen2-VL-2B-Instruct \\
Rank / $\alpha$ / dropout & 8 / 16 / 0.1 \\
Target modules & q, k, v, o, gate, up, down \\
Optimiser & AdamW \\
LR / weight decay & $5\times10^{-5}$ / 0.05 \\
Gradient accumulation & 4 steps \\
Epochs / examples & 1 / 16 \\
Hardware & GTX 1080 Ti (Ada cluster) \\
\bottomrule
\end{tabular}
\caption{LoRA configuration.}
\label{tab:ftconfig}
\end{table}

\subsection{Prompt and Output Parsing}

\paragraph{Classification.} Qwen2-VL-2B and LLaVA-Interleave-0.5B
received the two images, the full evidence card and the instruction
``Reply with exactly one word after the images: decline or increase or
stable. Do not copy the evidence text.'' PaliGemma received the two
images, the card's NDVI line and ``Two forest RGB chips: first image T0,
second T1. Vegetation change class --- answer with exactly one word:
stable, decline, or increase.'' BLIP-2 received the stitched image and
``Side by side: left forest patch T0, right forest patch T1. One word
only --- decline, increase, or stable:''. The \texttt{evidence\_only}
variant removes every card line that names \texttt{candidate\_class}.

\paragraph{Explanations.} The free-form prompt tells the model that
image 1 shows the patch in January--February 2022 and image 2 in
January--February 2023, asks it to describe the visible change in one
sentence, decide whether vegetation is stable, declined or increased,
and justify the decision. Qwen2-VL-2B and LLaVA answer in JSON;
PaliGemma and BLIP-2 receive a one-sentence version of the request.
The structured prompt adds a numbered list of allowed facts to be copied verbatim, a line stating
the $\pm 0.02$ NDVI rule for the patch, and a one-sentence justification
limited to NDVI. An automatic check confirms that no explanation prompt
contains \texttt{candidate\_class}.

\paragraph{Parsing.} The parser lowercases the output, accepts a
single-word answer directly, and otherwise searches for the first class
word or synonym. It does not exclude text copied from the prompt, which
explains BLIP-2's echo artefact.

\paragraph{Teacher filtering.} For the explanation LoRA, a generation
was kept only if its parsed class equalled the label, it contained at
least two checkable claims, and every claim agreed with the measured
change. This filter kept 28 of 256 training patches.

\subsection{Additional Results}

With $n = 55$, intervals are wide: 31/55 (56.4\%) corresponds to
43.3--68.6\% and 37/55 (67.3\%) to 54.1--78.2\%. Per-class counts
(about 18 patches per class) are too small to support a detailed error
taxonomy. The closed-set LLaVA probe (60.0\%) relies on a
development-fitted bias and is reported as a probe, not as a zero-shot
result; the same probe left PaliGemma and BLIP-2 at 32.7\%, predicting
``stable'' for every patch. In the \texttt{evidence\_only} runs, LLaVA
answered ``decline'' for 51 patches and ``increase'' for 4.
